% ── Chapter 5: Conclusion and Future Work ────────────────────
% Do NOT add \documentclass, \usepackage, \begin{document}, etc.
% This file is included by main.tex via % ── Chapter 5: Conclusion and Future Work ────────────────────
% Do NOT add \documentclass, \usepackage, \begin{document}, etc.
% This file is included by main.tex via % ── Chapter 5: Conclusion and Future Work ────────────────────
% Do NOT add \documentclass, \usepackage, \begin{document}, etc.
% This file is included by main.tex via % ── Chapter 5: Conclusion and Future Work ────────────────────
% Do NOT add \documentclass, \usepackage, \begin{document}, etc.
% This file is included by main.tex via \include{Chapter5}

\chapter{Conclusion and Future Work}

This dissertation presented the design, development, and experimental
evaluation of a communication-integrated mobile environmental monitoring
framework implemented on a Raspberry Pi~5 and Raspberry Pi~3 embedded
computing platform pair. The proposed framework addressed the fundamental
limitation of conventional fixed-station environmental monitoring by
integrating an autonomous mobile robotic platform with real-time
environmental sensing and multi-channel cloud-enabled data communication
within a modular and scalable architecture.

The framework was realised through four tightly integrated components.
The web-based manual control system provided a validated hardware
foundation through remote teleoperation, confirming the correct
operation of all mechanical, sensing, and communication subsystems.
The YOLOv8n-seg path segmentation module, trained on a custom
289-image annotated dataset using the CVAT annotation platform,
achieved a mask mAP@50 of 0.988 and a mask mAP@50--95 of 0.736 on
the 57-image validation split. Systematic optimisation through ONNX
format conversion and 320$\times$320 resolution tuning achieved
real-time inference at 15.4\,FPS with 64.8\,ms latency on the
Raspberry Pi~5 CPU, representing a 1.57$\times$ speedup over the
PyTorch baseline and exceeding the minimum frame rate required for
stable closed-loop control.

The autonomous path-following subsystem, driven by a PID controller
with gains $K_p = 13.2$, $K_i = 0.0$, and $K_d = 0.475$, operated
through a four-thread concurrent architecture at execution rates of
20\,Hz for vision inference, 10\,Hz for obstacle monitoring, and
10\,Hz for motor command execution. A systematic experimental
evaluation comprising 20 manual and 20 autonomous trials on a
rectangular 152.4\,cm $\times$ 91.44\,cm test path demonstrated that
autonomous navigation achieved a mean RMSE of 2.119\,cm, representing
an 87.4\% reduction from the manual mean RMSE of 16.768\,cm, and a
tracking accuracy of 95.37\% compared to 63.32\% under manual
teleoperation. All three statistical hypothesis tests (Mann--Whitney U,
permutation, and independent $t$-test) unanimously confirmed the
improvement to be highly significant, with an exceptionally large
Cohen's $d = 8.944$ and a probability of superiority of 100\%.

The environmental monitoring subsystem acquired temperature, humidity,
and air quality data from an Arduino-interfaced DHT11 and MQ135 sensor
pair, applied linear calibration to estimate CO$_2$ concentration in
ppm, and implemented a six-level rule-based alert hierarchy aligned
with ASHRAE Standard~62.1 and WHO Indoor Air Quality Guidelines. All
acquired readings were simultaneously persisted to a local SQLite
database, pushed to a Firebase Realtime Database hosted at the
asia-southeast1 region, transmitted to a ThingSpeak IoT channel with
six mapped sensor fields, and served through a Flask live web dashboard
with gauge displays, trend charts, and a reading history table. All
four parallel output channels operated concurrently without mutual
blocking, confirming the effectiveness of the threading architecture.

The integrated two-processor architecture, with the Raspberry Pi~5
managing autonomous navigation and the Raspberry Pi~3 managing
environmental data communication, demonstrated effective isolation
between the real-time navigation workload and the cloud communication
workload. Both subsystems operated concurrently on the shared mobile
platform without observable interference. The framework contributes a
practical and validated approach to communication-integrated mobile
sensing that combines deep learning perception, real-time embedded
control, and multi-platform IoT data communication on affordable
embedded hardware. The modular design supports future extensions,
establishing a replicable foundation for advanced mobile environmental
monitoring and smart infrastructure applications.

% ── Future Work ───────────────────────────────────────────────
\section{Future Work}

Several directions for future research and development are identified
based on the limitations encountered during implementation and the
results of experimental evaluation.

\subsection{Real Wheel Odometry Integration}

The current trajectory logging and position estimation system uses
simulated dead-reckoning from vision deviation and speed commands,
which accumulates drift over time. Integration of physical wheel
encoders or odometry sensors would enable accurate measurement of
distance travelled, angular velocity, and absolute position, replacing
the simulated estimate with ground-truth trajectory data and
significantly improving performance evaluation reliability.

\subsection{Model Quantisation and Hardware Acceleration}

The current ONNX model at 320$\times$320 resolution achieves 15.4\,FPS
on the Raspberry Pi~5 CPU. INT8 quantisation of the ONNX model is
expected to provide a further 2$\times$ speedup. Hardware acceleration
using the Coral USB Accelerator or Hailo AI HAT for the Raspberry Pi~5
could reach 25 to 30\,FPS, enabling smoother control at higher robot
speeds without requiring a full GPU platform.

\subsection{Extended Environmental Sensing}

The current subsystem measures temperature, humidity, and air quality
through the DHT11 and MQ135 sensors. Future work could extend the
sensing suite to include a dedicated CO$_2$ sensor such as the MH-Z19
for direct calibrated measurement, a particulate matter sensor such as
the PMS5003 for PM2.5 and PM10 monitoring, a volatile organic compound
sensor, and noise level monitoring, providing a more comprehensive
environmental data profile for smart building and occupational health
applications.

\subsection{Multi-Robot Distributed Framework}

The current framework operates as a single mobile unit. A natural
extension is the deployment of multiple mobile monitoring units that
collaborate to cover larger spatial areas. Each unit would transmit
environmental data to a shared cloud backend, enabling spatial mapping
of environmental parameters. Distributed task allocation algorithms
could coordinate path assignments between units to maximise coverage
efficiency.

\subsection{Advanced Communication Integration}

Future work could replace the current Wi-Fi-based HTTP communication
with cellular LTE or 5G modules to enable monitoring in areas without
fixed network infrastructure. Integration of LTE-V2X communication
frameworks could further extend the system toward vehicular
environmental monitoring applications in outdoor and urban environments.

% ── References ───────────────────────────────────────────────
\newpage
\renewcommand{\bibname}{References}
\addcontentsline{toc}{chapter}{References}

\begin{thebibliography}{50}

\bibitem{r1}
R. Pushkaranand and M. Reddy,
``IoT-Based Vehicular Pollution Monitoring System,''
\textit{IEEE International Conference on Green Computing and Internet of Things (ICGCIoT)},
2015.

\bibitem{r2}
Y. Ben-Aboud, M. Ghogho, and A. Kobbane,
``A Research-Oriented Low-Cost Air Pollution Monitoring IoT Platform,''
\textit{IEEE Conference},
2020.

\bibitem{r3}
P. Papadimitratos et al.,
``Secure Vehicular Communication Systems: Design and Architecture,''
\textit{IEEE Communications Magazine},
vol.~46, no.~11, pp.~100--109, Nov. 2008.

\bibitem{r4}
G. Araniti, C. Campolo, M. Condoluci, A. Iera, and A. Molinaro,
``LTE for Vehicular Networking: A Survey,''
\textit{IEEE Communications Magazine},
vol.~51, no.~5, pp.~148--157, May 2013.

\bibitem{r5}
Y. Ma, J. Ko\v{s}eck\'{a}, and S. Sastry,
``Vision Guided Navigation for a Nonholonomic Mobile Robot,''
\textit{Proceedings of the IEEE International Conference on Robotics and Automation (ICRA)},
1999.

\bibitem{r6}
Y. Fang, W. Dixon, D. Dawson, and C. Abdallah,
``Homography-Based Visual Servoing of Wheeled Mobile Robots,''
\textit{IEEE Transactions on Robotics},
vol.~21, no.~3, pp.~406--414, 2005.

\bibitem{r7}
Y. Fang, W. Dixon, D. Dawson, and C. Abdallah,
``Homography-Based Visual Servo Regulation for Mobile Robots,''
\textit{IEEE Transactions on Robotics},
vol.~21, no.~3, pp.~406--414, 2005.

\bibitem{r8}
J. Chen, W. Dixon, D. Dawson, and M. McIntyre,
``Homography-Based Visual Servo Tracking Control of a Wheeled Mobile Robot,''
\textit{IEEE Transactions on Robotics},
2006.

\bibitem{r9}
E. Shaban and A. Nada,
``PID vs PIP Control for Mobile Robots,''
\textit{International Journal of Computer Applications},
vol.~117, no.~2, 2015.

\bibitem{r10}
L. Pacheco and N. Luo,
``PID vs MPC Control Strategies for Mobile Robots: A Comparative Study,''
\textit{IEEE Conference on Control Applications},
2015.

\bibitem{r11}
N. Noordin et al.,
``Comparison of PID and Fuzzy Logic Control for Mobile Robots,''
\textit{International Journal of Robotics Research},
2021.

\bibitem{r12}
M. Siam, S. Elkerdawy, M. Jagersand, and S. Yogamani,
``Deep Semantic Segmentation for Automated Driving: A Survey,''
\textit{IEEE Transactions on Intelligent Transportation Systems},
vol.~20, no.~2, pp.~641--659, 2018.

\bibitem{r13}
J. Dos Reis et al.,
``YOLO-Based Object Detection for Autonomous Navigation,''
\textit{IEEE Access},
vol.~7, pp.~124908--124920, 2019.

\bibitem{r14}
J. Jean and K. Lian,
``Robust Visual Servo Control Using Shape Parameters,''
\textit{IEEE/ASME Transactions on Mechatronics},
vol.~17, no.~1, pp.~30--40, 2012.

\bibitem{r15}
A. Zde\v{s}ar, D. Matko, and M. Bla\v{z}i\v{c},
``Visual Trajectory Tracking Control of Mobile Robots,''
\textit{IEEE Transactions on Control Systems Technology},
vol.~21, no.~6, pp.~2210--2217, 2013.

\bibitem{r16}
S. Mir and F. Filali,
``LTE and IEEE 802.11p for Vehicular Networking: A Performance Evaluation,''
\textit{EURASIP Journal on Wireless Communications and Networking},
2014.

\bibitem{r17}
K. Yousif, A. Bab-Hadiashar, and R. Hoseinnezhad,
``An Overview to Visual Odometry and Visual SLAM,''
\textit{Applied Bionics and Biomechanics},
2015.

\bibitem{r18}
C. Pati and R. Kala,
``Vision-Based Robot Following Using RGB-D Sensors,''
\textit{Procedia Computer Science},
vol.~115, pp.~586--593, 2017.

\bibitem{r19}
Z. Xu et al.,
``DSRC versus 4G-LTE for Connected Vehicle Applications: A Study on Field Experiments,''
\textit{Journal of Advanced Transportation},
2017.

\bibitem{r20}
Y. Zhang, Z. Zhu, and H. Wang,
``Global Homography Calibration for Vision-Based Robot Navigation,''
\textit{IEEE Transactions on Instrumentation and Measurement},
vol.~66, no.~10, pp.~2754--2763, 2017.

\bibitem{r21}
M. Fauadi et al.,
``A Review on Vision-Based Navigation Systems for Mobile Robots,''
\textit{International Journal of Advanced Robotic Systems},
2018.

\bibitem{r22}
J. M. Lozano Dom\'{i}nguez and T. J. Mateo Sanguino,
``Review on V2X, I2X, and P2X Communications and Their Applications,''
\textit{Sensors},
vol.~19, no.~12, p.~2756, 2019.

\bibitem{r23}
P. S. Santana et al.,
``Vehicle-Mounted Sensing for Mobile Air Quality Monitoring,''
\textit{IEEE Conference on Intelligent Transportation Systems},
2019.

\bibitem{r24}
P. Papadeas, G. Athanasopoulos, and K. Karantzalos,
``A Survey on Real-Time Semantic Segmentation for Autonomous Systems,''
\textit{IEEE Access},
vol.~9, pp.~54215--54240, 2021.

\bibitem{r25}
T. Petrov et al.,
``Performance Benchmark for DSRC and LTE-V2X Communication,''
\textit{Sensors},
vol.~21, no.~6, 2021.

\bibitem{r26}
A. Singh, R. Sharma, and P. Kumar,
``Machine Learning-Based Obstacle Avoidance and Path Tracking for Mobile Robots,''
\textit{IEEE Access},
vol.~10, pp.~112345--112356, 2022.

\bibitem{r27}
A. Saidi, H. Mekki, and M. Abid,
``Fuzzy-PID Control for Vision-Based Mobile Robot Navigation,''
\textit{Robotics and Autonomous Systems},
vol.~150, 2022.

\bibitem{r28}
C. Shin, E. Farag, H. Ryu, M. Zhou, and Y. Kim,
``Vehicle-to-Everything (V2X) Evolution from 4G to 5G in 3GPP,''
\textit{IEEE Access},
vol.~11, pp.~18689--18701, 2023.

\bibitem{r29}
J. Marroqu\'{i}n, R. Monroy, and L. Castillo,
``Embedded Computer Vision Navigation System for Mobile Robots,''
\textit{IEEE Access},
vol.~11, pp.~56789--56805, 2023.

\bibitem{r30}
M. M\i s\i r,
``DeepLab-Based Drivable Path Detection for Mobile Robot Navigation,''
\textit{IEEE Access},
vol.~12, pp.~12345--12360, 2024.

\bibitem{r31}
CVAT.ai Corporation,
\textit{CVAT: Computer Vision Annotation Tool},
2024.
[Online]. Available: \url{https://github.com/cvat-ai/cvat}

\bibitem{r32}
Ultralytics,
\textit{YOLOv8: Ultralytics YOLO Documentation},
Version 8.4.7, 2024.
[Online]. Available: \url{https://docs.ultralytics.com}

\bibitem{r33}
ONNX Runtime Development Team,
\textit{ONNX Runtime: Cross-Platform Inference Accelerator},
Version 1.23.2, 2024.
[Online]. Available: \url{https://onnxruntime.ai/docs/}

\bibitem{r34}
Raspberry Pi Foundation,
\textit{Raspberry Pi 5 Official Documentation},
2024.
[Online]. Available: \url{https://www.raspberrypi.com/documentation/}

\bibitem{r35}
Google Firebase,
\textit{Firebase Realtime Database Documentation},
Google LLC, 2024.
[Online]. Available: \url{https://firebase.google.com/docs/database}

\bibitem{r36}
MathWorks ThingSpeak,
\textit{ThingSpeak IoT Analytics Platform Documentation},
MathWorks, 2024.
[Online]. Available: \url{https://uk.mathworks.com/help/thingspeak/}

\bibitem{r37}
Pallets Projects,
\textit{Flask: A Lightweight WSGI Web Application Framework},
Version 3.x, 2024.
[Online]. Available: \url{https://flask.palletsprojects.com/}

\bibitem{r38}
Adafruit Industries,
\textit{Adafruit CircuitPython PCA9685 Library},
2024.
[Online]. Available: \url{https://github.com/adafruit/Adafruit_CircuitPython_PCA9685}

\bibitem{r39}
F. Pedregosa et al.,
``Scikit-learn: Machine Learning in Python,''
\textit{Journal of Machine Learning Research},
vol.~12, pp.~2825--2830, 2011.
[Online]. Available: \url{https://scikit-learn.org}

\bibitem{r40}
C. R. Harris et al.,
``Array Programming with NumPy,''
\textit{Nature},
vol.~585, pp.~357--362, 2020.
[Online]. Available: \url{https://numpy.org}

\bibitem{r41}
P. Virtanen et al.,
``SciPy 1.0: Fundamental Algorithms for Scientific Computing in Python,''
\textit{Nature Methods},
vol.~17, pp.~261--272, 2020.
[Online]. Available: \url{https://scipy.org}

\bibitem{r42}
J. D. Hunter,
``Matplotlib: A 2D Graphics Environment,''
\textit{Computing in Science \& Engineering},
vol.~9, no.~3, pp.~90--95, 2007.
[Online]. Available: \url{https://matplotlib.org}

\bibitem{r43}
OpenCV Development Team,
\textit{OpenCV: Open Source Computer Vision Library},
Version 4.x, 2024.
[Online]. Available: \url{https://opencv.org}

\bibitem{r44}
PySerial Development Team,
\textit{PySerial: Python Serial Port Extension},
Version 3.x, 2024.
[Online]. Available: \url{https://github.com/pyserial/pyserial}

\bibitem{r45}
Google Cloud,
\textit{Firebase Admin Python SDK},
Version 6.x, 2024.
[Online]. Available: \url{https://github.com/firebase/firebase-admin-python}

\bibitem{r46}
American Society of Heating, Refrigerating and Air-Conditioning Engineers (ASHRAE),
\textit{ASHRAE Standard 62.1: Ventilation for Acceptable Indoor Air Quality},
ASHRAE, Atlanta, GA, 2022.

\bibitem{r47}
World Health Organization (WHO),
\textit{WHO Global Air Quality Guidelines},
WHO, Geneva, 2021.
[Online]. Available: \url{https://www.who.int/publications/i/item/9789240034228}

\bibitem{r48}
O.\ Sepp\"{a}nen, W.\,J.\ Fisk, and M.\,J.\ Mendell,
``Association of Ventilation Rates and CO$_2$ Concentrations with Health
and Other Responses in Commercial and Institutional Buildings,''
\textit{Indoor Air},
vol.~9, no.~4, pp.~226--252, 1999.

\bibitem{r49}
U.\ Satish et al.,
``Is CO$_2$ an Indoor Pollutant? Direct Effects of Low-to-Moderate
CO$_2$ Concentrations on Human Decision-Making Performance,''
\textit{Environmental Health Perspectives},
vol.~120, no.~12, pp.~1671--1677, 2012.

\bibitem{r50}
Hanwei Electronics Co.,
\textit{MQ-135 Gas Sensor for Air Quality Datasheet},
Hanwei Electronics, Zhengzhou, China, 2013.
[Online]. Available: \url{https://www.hwsensor.com}

\end{thebibliography}

\chapter{Conclusion and Future Work}

This dissertation presented the design, development, and experimental
evaluation of a communication-integrated mobile environmental monitoring
framework implemented on a Raspberry Pi~5 and Raspberry Pi~3 embedded
computing platform pair. The proposed framework addressed the fundamental
limitation of conventional fixed-station environmental monitoring by
integrating an autonomous mobile robotic platform with real-time
environmental sensing and multi-channel cloud-enabled data communication
within a modular and scalable architecture.

The framework was realised through four tightly integrated components.
The web-based manual control system provided a validated hardware
foundation through remote teleoperation, confirming the correct
operation of all mechanical, sensing, and communication subsystems.
The YOLOv8n-seg path segmentation module, trained on a custom
289-image annotated dataset using the CVAT annotation platform,
achieved a mask mAP@50 of 0.988 and a mask mAP@50--95 of 0.736 on
the 57-image validation split. Systematic optimisation through ONNX
format conversion and 320$\times$320 resolution tuning achieved
real-time inference at 15.4\,FPS with 64.8\,ms latency on the
Raspberry Pi~5 CPU, representing a 1.57$\times$ speedup over the
PyTorch baseline and exceeding the minimum frame rate required for
stable closed-loop control.

The autonomous path-following subsystem, driven by a PID controller
with gains $K_p = 13.2$, $K_i = 0.0$, and $K_d = 0.475$, operated
through a four-thread concurrent architecture at execution rates of
20\,Hz for vision inference, 10\,Hz for obstacle monitoring, and
10\,Hz for motor command execution. A systematic experimental
evaluation comprising 20 manual and 20 autonomous trials on a
rectangular 152.4\,cm $\times$ 91.44\,cm test path demonstrated that
autonomous navigation achieved a mean RMSE of 2.119\,cm, representing
an 87.4\% reduction from the manual mean RMSE of 16.768\,cm, and a
tracking accuracy of 95.37\% compared to 63.32\% under manual
teleoperation. All three statistical hypothesis tests (Mann--Whitney U,
permutation, and independent $t$-test) unanimously confirmed the
improvement to be highly significant, with an exceptionally large
Cohen's $d = 8.944$ and a probability of superiority of 100\%.

The environmental monitoring subsystem acquired temperature, humidity,
and air quality data from an Arduino-interfaced DHT11 and MQ135 sensor
pair, applied linear calibration to estimate CO$_2$ concentration in
ppm, and implemented a six-level rule-based alert hierarchy aligned
with ASHRAE Standard~62.1 and WHO Indoor Air Quality Guidelines. All
acquired readings were simultaneously persisted to a local SQLite
database, pushed to a Firebase Realtime Database hosted at the
asia-southeast1 region, transmitted to a ThingSpeak IoT channel with
six mapped sensor fields, and served through a Flask live web dashboard
with gauge displays, trend charts, and a reading history table. All
four parallel output channels operated concurrently without mutual
blocking, confirming the effectiveness of the threading architecture.

The integrated two-processor architecture, with the Raspberry Pi~5
managing autonomous navigation and the Raspberry Pi~3 managing
environmental data communication, demonstrated effective isolation
between the real-time navigation workload and the cloud communication
workload. Both subsystems operated concurrently on the shared mobile
platform without observable interference. The framework contributes a
practical and validated approach to communication-integrated mobile
sensing that combines deep learning perception, real-time embedded
control, and multi-platform IoT data communication on affordable
embedded hardware. The modular design supports future extensions,
establishing a replicable foundation for advanced mobile environmental
monitoring and smart infrastructure applications.

% ── Future Work ───────────────────────────────────────────────
\section{Future Work}

Several directions for future research and development are identified
based on the limitations encountered during implementation and the
results of experimental evaluation.

\subsection{Real Wheel Odometry Integration}

The current trajectory logging and position estimation system uses
simulated dead-reckoning from vision deviation and speed commands,
which accumulates drift over time. Integration of physical wheel
encoders or odometry sensors would enable accurate measurement of
distance travelled, angular velocity, and absolute position, replacing
the simulated estimate with ground-truth trajectory data and
significantly improving performance evaluation reliability.

\subsection{Model Quantisation and Hardware Acceleration}

The current ONNX model at 320$\times$320 resolution achieves 15.4\,FPS
on the Raspberry Pi~5 CPU. INT8 quantisation of the ONNX model is
expected to provide a further 2$\times$ speedup. Hardware acceleration
using the Coral USB Accelerator or Hailo AI HAT for the Raspberry Pi~5
could reach 25 to 30\,FPS, enabling smoother control at higher robot
speeds without requiring a full GPU platform.

\subsection{Extended Environmental Sensing}

The current subsystem measures temperature, humidity, and air quality
through the DHT11 and MQ135 sensors. Future work could extend the
sensing suite to include a dedicated CO$_2$ sensor such as the MH-Z19
for direct calibrated measurement, a particulate matter sensor such as
the PMS5003 for PM2.5 and PM10 monitoring, a volatile organic compound
sensor, and noise level monitoring, providing a more comprehensive
environmental data profile for smart building and occupational health
applications.

\subsection{Multi-Robot Distributed Framework}

The current framework operates as a single mobile unit. A natural
extension is the deployment of multiple mobile monitoring units that
collaborate to cover larger spatial areas. Each unit would transmit
environmental data to a shared cloud backend, enabling spatial mapping
of environmental parameters. Distributed task allocation algorithms
could coordinate path assignments between units to maximise coverage
efficiency.

\subsection{Advanced Communication Integration}

Future work could replace the current Wi-Fi-based HTTP communication
with cellular LTE or 5G modules to enable monitoring in areas without
fixed network infrastructure. Integration of LTE-V2X communication
frameworks could further extend the system toward vehicular
environmental monitoring applications in outdoor and urban environments.

% ── References ───────────────────────────────────────────────
\newpage
\renewcommand{\bibname}{References}
\addcontentsline{toc}{chapter}{References}

\begin{thebibliography}{50}

\bibitem{r1}
R. Pushkaranand and M. Reddy,
``IoT-Based Vehicular Pollution Monitoring System,''
\textit{IEEE International Conference on Green Computing and Internet of Things (ICGCIoT)},
2015.

\bibitem{r2}
Y. Ben-Aboud, M. Ghogho, and A. Kobbane,
``A Research-Oriented Low-Cost Air Pollution Monitoring IoT Platform,''
\textit{IEEE Conference},
2020.

\bibitem{r3}
P. Papadimitratos et al.,
``Secure Vehicular Communication Systems: Design and Architecture,''
\textit{IEEE Communications Magazine},
vol.~46, no.~11, pp.~100--109, Nov. 2008.

\bibitem{r4}
G. Araniti, C. Campolo, M. Condoluci, A. Iera, and A. Molinaro,
``LTE for Vehicular Networking: A Survey,''
\textit{IEEE Communications Magazine},
vol.~51, no.~5, pp.~148--157, May 2013.

\bibitem{r5}
Y. Ma, J. Ko\v{s}eck\'{a}, and S. Sastry,
``Vision Guided Navigation for a Nonholonomic Mobile Robot,''
\textit{Proceedings of the IEEE International Conference on Robotics and Automation (ICRA)},
1999.

\bibitem{r6}
Y. Fang, W. Dixon, D. Dawson, and C. Abdallah,
``Homography-Based Visual Servoing of Wheeled Mobile Robots,''
\textit{IEEE Transactions on Robotics},
vol.~21, no.~3, pp.~406--414, 2005.

\bibitem{r7}
Y. Fang, W. Dixon, D. Dawson, and C. Abdallah,
``Homography-Based Visual Servo Regulation for Mobile Robots,''
\textit{IEEE Transactions on Robotics},
vol.~21, no.~3, pp.~406--414, 2005.

\bibitem{r8}
J. Chen, W. Dixon, D. Dawson, and M. McIntyre,
``Homography-Based Visual Servo Tracking Control of a Wheeled Mobile Robot,''
\textit{IEEE Transactions on Robotics},
2006.

\bibitem{r9}
E. Shaban and A. Nada,
``PID vs PIP Control for Mobile Robots,''
\textit{International Journal of Computer Applications},
vol.~117, no.~2, 2015.

\bibitem{r10}
L. Pacheco and N. Luo,
``PID vs MPC Control Strategies for Mobile Robots: A Comparative Study,''
\textit{IEEE Conference on Control Applications},
2015.

\bibitem{r11}
N. Noordin et al.,
``Comparison of PID and Fuzzy Logic Control for Mobile Robots,''
\textit{International Journal of Robotics Research},
2021.

\bibitem{r12}
M. Siam, S. Elkerdawy, M. Jagersand, and S. Yogamani,
``Deep Semantic Segmentation for Automated Driving: A Survey,''
\textit{IEEE Transactions on Intelligent Transportation Systems},
vol.~20, no.~2, pp.~641--659, 2018.

\bibitem{r13}
J. Dos Reis et al.,
``YOLO-Based Object Detection for Autonomous Navigation,''
\textit{IEEE Access},
vol.~7, pp.~124908--124920, 2019.

\bibitem{r14}
J. Jean and K. Lian,
``Robust Visual Servo Control Using Shape Parameters,''
\textit{IEEE/ASME Transactions on Mechatronics},
vol.~17, no.~1, pp.~30--40, 2012.

\bibitem{r15}
A. Zde\v{s}ar, D. Matko, and M. Bla\v{z}i\v{c},
``Visual Trajectory Tracking Control of Mobile Robots,''
\textit{IEEE Transactions on Control Systems Technology},
vol.~21, no.~6, pp.~2210--2217, 2013.

\bibitem{r16}
S. Mir and F. Filali,
``LTE and IEEE 802.11p for Vehicular Networking: A Performance Evaluation,''
\textit{EURASIP Journal on Wireless Communications and Networking},
2014.

\bibitem{r17}
K. Yousif, A. Bab-Hadiashar, and R. Hoseinnezhad,
``An Overview to Visual Odometry and Visual SLAM,''
\textit{Applied Bionics and Biomechanics},
2015.

\bibitem{r18}
C. Pati and R. Kala,
``Vision-Based Robot Following Using RGB-D Sensors,''
\textit{Procedia Computer Science},
vol.~115, pp.~586--593, 2017.

\bibitem{r19}
Z. Xu et al.,
``DSRC versus 4G-LTE for Connected Vehicle Applications: A Study on Field Experiments,''
\textit{Journal of Advanced Transportation},
2017.

\bibitem{r20}
Y. Zhang, Z. Zhu, and H. Wang,
``Global Homography Calibration for Vision-Based Robot Navigation,''
\textit{IEEE Transactions on Instrumentation and Measurement},
vol.~66, no.~10, pp.~2754--2763, 2017.

\bibitem{r21}
M. Fauadi et al.,
``A Review on Vision-Based Navigation Systems for Mobile Robots,''
\textit{International Journal of Advanced Robotic Systems},
2018.

\bibitem{r22}
J. M. Lozano Dom\'{i}nguez and T. J. Mateo Sanguino,
``Review on V2X, I2X, and P2X Communications and Their Applications,''
\textit{Sensors},
vol.~19, no.~12, p.~2756, 2019.

\bibitem{r23}
P. S. Santana et al.,
``Vehicle-Mounted Sensing for Mobile Air Quality Monitoring,''
\textit{IEEE Conference on Intelligent Transportation Systems},
2019.

\bibitem{r24}
P. Papadeas, G. Athanasopoulos, and K. Karantzalos,
``A Survey on Real-Time Semantic Segmentation for Autonomous Systems,''
\textit{IEEE Access},
vol.~9, pp.~54215--54240, 2021.

\bibitem{r25}
T. Petrov et al.,
``Performance Benchmark for DSRC and LTE-V2X Communication,''
\textit{Sensors},
vol.~21, no.~6, 2021.

\bibitem{r26}
A. Singh, R. Sharma, and P. Kumar,
``Machine Learning-Based Obstacle Avoidance and Path Tracking for Mobile Robots,''
\textit{IEEE Access},
vol.~10, pp.~112345--112356, 2022.

\bibitem{r27}
A. Saidi, H. Mekki, and M. Abid,
``Fuzzy-PID Control for Vision-Based Mobile Robot Navigation,''
\textit{Robotics and Autonomous Systems},
vol.~150, 2022.

\bibitem{r28}
C. Shin, E. Farag, H. Ryu, M. Zhou, and Y. Kim,
``Vehicle-to-Everything (V2X) Evolution from 4G to 5G in 3GPP,''
\textit{IEEE Access},
vol.~11, pp.~18689--18701, 2023.

\bibitem{r29}
J. Marroqu\'{i}n, R. Monroy, and L. Castillo,
``Embedded Computer Vision Navigation System for Mobile Robots,''
\textit{IEEE Access},
vol.~11, pp.~56789--56805, 2023.

\bibitem{r30}
M. M\i s\i r,
``DeepLab-Based Drivable Path Detection for Mobile Robot Navigation,''
\textit{IEEE Access},
vol.~12, pp.~12345--12360, 2024.

\bibitem{r31}
CVAT.ai Corporation,
\textit{CVAT: Computer Vision Annotation Tool},
2024.
[Online]. Available: \url{https://github.com/cvat-ai/cvat}

\bibitem{r32}
Ultralytics,
\textit{YOLOv8: Ultralytics YOLO Documentation},
Version 8.4.7, 2024.
[Online]. Available: \url{https://docs.ultralytics.com}

\bibitem{r33}
ONNX Runtime Development Team,
\textit{ONNX Runtime: Cross-Platform Inference Accelerator},
Version 1.23.2, 2024.
[Online]. Available: \url{https://onnxruntime.ai/docs/}

\bibitem{r34}
Raspberry Pi Foundation,
\textit{Raspberry Pi 5 Official Documentation},
2024.
[Online]. Available: \url{https://www.raspberrypi.com/documentation/}

\bibitem{r35}
Google Firebase,
\textit{Firebase Realtime Database Documentation},
Google LLC, 2024.
[Online]. Available: \url{https://firebase.google.com/docs/database}

\bibitem{r36}
MathWorks ThingSpeak,
\textit{ThingSpeak IoT Analytics Platform Documentation},
MathWorks, 2024.
[Online]. Available: \url{https://uk.mathworks.com/help/thingspeak/}

\bibitem{r37}
Pallets Projects,
\textit{Flask: A Lightweight WSGI Web Application Framework},
Version 3.x, 2024.
[Online]. Available: \url{https://flask.palletsprojects.com/}

\bibitem{r38}
Adafruit Industries,
\textit{Adafruit CircuitPython PCA9685 Library},
2024.
[Online]. Available: \url{https://github.com/adafruit/Adafruit_CircuitPython_PCA9685}

\bibitem{r39}
F. Pedregosa et al.,
``Scikit-learn: Machine Learning in Python,''
\textit{Journal of Machine Learning Research},
vol.~12, pp.~2825--2830, 2011.
[Online]. Available: \url{https://scikit-learn.org}

\bibitem{r40}
C. R. Harris et al.,
``Array Programming with NumPy,''
\textit{Nature},
vol.~585, pp.~357--362, 2020.
[Online]. Available: \url{https://numpy.org}

\bibitem{r41}
P. Virtanen et al.,
``SciPy 1.0: Fundamental Algorithms for Scientific Computing in Python,''
\textit{Nature Methods},
vol.~17, pp.~261--272, 2020.
[Online]. Available: \url{https://scipy.org}

\bibitem{r42}
J. D. Hunter,
``Matplotlib: A 2D Graphics Environment,''
\textit{Computing in Science \& Engineering},
vol.~9, no.~3, pp.~90--95, 2007.
[Online]. Available: \url{https://matplotlib.org}

\bibitem{r43}
OpenCV Development Team,
\textit{OpenCV: Open Source Computer Vision Library},
Version 4.x, 2024.
[Online]. Available: \url{https://opencv.org}

\bibitem{r44}
PySerial Development Team,
\textit{PySerial: Python Serial Port Extension},
Version 3.x, 2024.
[Online]. Available: \url{https://github.com/pyserial/pyserial}

\bibitem{r45}
Google Cloud,
\textit{Firebase Admin Python SDK},
Version 6.x, 2024.
[Online]. Available: \url{https://github.com/firebase/firebase-admin-python}

\bibitem{r46}
American Society of Heating, Refrigerating and Air-Conditioning Engineers (ASHRAE),
\textit{ASHRAE Standard 62.1: Ventilation for Acceptable Indoor Air Quality},
ASHRAE, Atlanta, GA, 2022.

\bibitem{r47}
World Health Organization (WHO),
\textit{WHO Global Air Quality Guidelines},
WHO, Geneva, 2021.
[Online]. Available: \url{https://www.who.int/publications/i/item/9789240034228}

\bibitem{r48}
O.\ Sepp\"{a}nen, W.\,J.\ Fisk, and M.\,J.\ Mendell,
``Association of Ventilation Rates and CO$_2$ Concentrations with Health
and Other Responses in Commercial and Institutional Buildings,''
\textit{Indoor Air},
vol.~9, no.~4, pp.~226--252, 1999.

\bibitem{r49}
U.\ Satish et al.,
``Is CO$_2$ an Indoor Pollutant? Direct Effects of Low-to-Moderate
CO$_2$ Concentrations on Human Decision-Making Performance,''
\textit{Environmental Health Perspectives},
vol.~120, no.~12, pp.~1671--1677, 2012.

\bibitem{r50}
Hanwei Electronics Co.,
\textit{MQ-135 Gas Sensor for Air Quality Datasheet},
Hanwei Electronics, Zhengzhou, China, 2013.
[Online]. Available: \url{https://www.hwsensor.com}

\end{thebibliography}

\chapter{Conclusion and Future Work}

This dissertation presented the design, development, and experimental
evaluation of a communication-integrated mobile environmental monitoring
framework implemented on a Raspberry Pi~5 and Raspberry Pi~3 embedded
computing platform pair. The proposed framework addressed the fundamental
limitation of conventional fixed-station environmental monitoring by
integrating an autonomous mobile robotic platform with real-time
environmental sensing and multi-channel cloud-enabled data communication
within a modular and scalable architecture.

The framework was realised through four tightly integrated components.
The web-based manual control system provided a validated hardware
foundation through remote teleoperation, confirming the correct
operation of all mechanical, sensing, and communication subsystems.
The YOLOv8n-seg path segmentation module, trained on a custom
289-image annotated dataset using the CVAT annotation platform,
achieved a mask mAP@50 of 0.988 and a mask mAP@50--95 of 0.736 on
the 57-image validation split. Systematic optimisation through ONNX
format conversion and 320$\times$320 resolution tuning achieved
real-time inference at 15.4\,FPS with 64.8\,ms latency on the
Raspberry Pi~5 CPU, representing a 1.57$\times$ speedup over the
PyTorch baseline and exceeding the minimum frame rate required for
stable closed-loop control.

The autonomous path-following subsystem, driven by a PID controller
with gains $K_p = 13.2$, $K_i = 0.0$, and $K_d = 0.475$, operated
through a four-thread concurrent architecture at execution rates of
20\,Hz for vision inference, 10\,Hz for obstacle monitoring, and
10\,Hz for motor command execution. A systematic experimental
evaluation comprising 20 manual and 20 autonomous trials on a
rectangular 152.4\,cm $\times$ 91.44\,cm test path demonstrated that
autonomous navigation achieved a mean RMSE of 2.119\,cm, representing
an 87.4\% reduction from the manual mean RMSE of 16.768\,cm, and a
tracking accuracy of 95.37\% compared to 63.32\% under manual
teleoperation. All three statistical hypothesis tests (Mann--Whitney U,
permutation, and independent $t$-test) unanimously confirmed the
improvement to be highly significant, with an exceptionally large
Cohen's $d = 8.944$ and a probability of superiority of 100\%.

The environmental monitoring subsystem acquired temperature, humidity,
and air quality data from an Arduino-interfaced DHT11 and MQ135 sensor
pair, applied linear calibration to estimate CO$_2$ concentration in
ppm, and implemented a six-level rule-based alert hierarchy aligned
with ASHRAE Standard~62.1 and WHO Indoor Air Quality Guidelines. All
acquired readings were simultaneously persisted to a local SQLite
database, pushed to a Firebase Realtime Database hosted at the
asia-southeast1 region, transmitted to a ThingSpeak IoT channel with
six mapped sensor fields, and served through a Flask live web dashboard
with gauge displays, trend charts, and a reading history table. All
four parallel output channels operated concurrently without mutual
blocking, confirming the effectiveness of the threading architecture.

The integrated two-processor architecture, with the Raspberry Pi~5
managing autonomous navigation and the Raspberry Pi~3 managing
environmental data communication, demonstrated effective isolation
between the real-time navigation workload and the cloud communication
workload. Both subsystems operated concurrently on the shared mobile
platform without observable interference. The framework contributes a
practical and validated approach to communication-integrated mobile
sensing that combines deep learning perception, real-time embedded
control, and multi-platform IoT data communication on affordable
embedded hardware. The modular design supports future extensions,
establishing a replicable foundation for advanced mobile environmental
monitoring and smart infrastructure applications.

% ── Future Work ───────────────────────────────────────────────
\section{Future Work}

Several directions for future research and development are identified
based on the limitations encountered during implementation and the
results of experimental evaluation.

\subsection{Real Wheel Odometry Integration}

The current trajectory logging and position estimation system uses
simulated dead-reckoning from vision deviation and speed commands,
which accumulates drift over time. Integration of physical wheel
encoders or odometry sensors would enable accurate measurement of
distance travelled, angular velocity, and absolute position, replacing
the simulated estimate with ground-truth trajectory data and
significantly improving performance evaluation reliability.

\subsection{Model Quantisation and Hardware Acceleration}

The current ONNX model at 320$\times$320 resolution achieves 15.4\,FPS
on the Raspberry Pi~5 CPU. INT8 quantisation of the ONNX model is
expected to provide a further 2$\times$ speedup. Hardware acceleration
using the Coral USB Accelerator or Hailo AI HAT for the Raspberry Pi~5
could reach 25 to 30\,FPS, enabling smoother control at higher robot
speeds without requiring a full GPU platform.

\subsection{Extended Environmental Sensing}

The current subsystem measures temperature, humidity, and air quality
through the DHT11 and MQ135 sensors. Future work could extend the
sensing suite to include a dedicated CO$_2$ sensor such as the MH-Z19
for direct calibrated measurement, a particulate matter sensor such as
the PMS5003 for PM2.5 and PM10 monitoring, a volatile organic compound
sensor, and noise level monitoring, providing a more comprehensive
environmental data profile for smart building and occupational health
applications.

\subsection{Multi-Robot Distributed Framework}

The current framework operates as a single mobile unit. A natural
extension is the deployment of multiple mobile monitoring units that
collaborate to cover larger spatial areas. Each unit would transmit
environmental data to a shared cloud backend, enabling spatial mapping
of environmental parameters. Distributed task allocation algorithms
could coordinate path assignments between units to maximise coverage
efficiency.

\subsection{Advanced Communication Integration}

Future work could replace the current Wi-Fi-based HTTP communication
with cellular LTE or 5G modules to enable monitoring in areas without
fixed network infrastructure. Integration of LTE-V2X communication
frameworks could further extend the system toward vehicular
environmental monitoring applications in outdoor and urban environments.

% ── References ───────────────────────────────────────────────
\newpage
\renewcommand{\bibname}{References}
\addcontentsline{toc}{chapter}{References}

\begin{thebibliography}{50}

\bibitem{r1}
R. Pushkaranand and M. Reddy,
``IoT-Based Vehicular Pollution Monitoring System,''
\textit{IEEE International Conference on Green Computing and Internet of Things (ICGCIoT)},
2015.

\bibitem{r2}
Y. Ben-Aboud, M. Ghogho, and A. Kobbane,
``A Research-Oriented Low-Cost Air Pollution Monitoring IoT Platform,''
\textit{IEEE Conference},
2020.

\bibitem{r3}
P. Papadimitratos et al.,
``Secure Vehicular Communication Systems: Design and Architecture,''
\textit{IEEE Communications Magazine},
vol.~46, no.~11, pp.~100--109, Nov. 2008.

\bibitem{r4}
G. Araniti, C. Campolo, M. Condoluci, A. Iera, and A. Molinaro,
``LTE for Vehicular Networking: A Survey,''
\textit{IEEE Communications Magazine},
vol.~51, no.~5, pp.~148--157, May 2013.

\bibitem{r5}
Y. Ma, J. Ko\v{s}eck\'{a}, and S. Sastry,
``Vision Guided Navigation for a Nonholonomic Mobile Robot,''
\textit{Proceedings of the IEEE International Conference on Robotics and Automation (ICRA)},
1999.

\bibitem{r6}
Y. Fang, W. Dixon, D. Dawson, and C. Abdallah,
``Homography-Based Visual Servoing of Wheeled Mobile Robots,''
\textit{IEEE Transactions on Robotics},
vol.~21, no.~3, pp.~406--414, 2005.

\bibitem{r7}
Y. Fang, W. Dixon, D. Dawson, and C. Abdallah,
``Homography-Based Visual Servo Regulation for Mobile Robots,''
\textit{IEEE Transactions on Robotics},
vol.~21, no.~3, pp.~406--414, 2005.

\bibitem{r8}
J. Chen, W. Dixon, D. Dawson, and M. McIntyre,
``Homography-Based Visual Servo Tracking Control of a Wheeled Mobile Robot,''
\textit{IEEE Transactions on Robotics},
2006.

\bibitem{r9}
E. Shaban and A. Nada,
``PID vs PIP Control for Mobile Robots,''
\textit{International Journal of Computer Applications},
vol.~117, no.~2, 2015.

\bibitem{r10}
L. Pacheco and N. Luo,
``PID vs MPC Control Strategies for Mobile Robots: A Comparative Study,''
\textit{IEEE Conference on Control Applications},
2015.

\bibitem{r11}
N. Noordin et al.,
``Comparison of PID and Fuzzy Logic Control for Mobile Robots,''
\textit{International Journal of Robotics Research},
2021.

\bibitem{r12}
M. Siam, S. Elkerdawy, M. Jagersand, and S. Yogamani,
``Deep Semantic Segmentation for Automated Driving: A Survey,''
\textit{IEEE Transactions on Intelligent Transportation Systems},
vol.~20, no.~2, pp.~641--659, 2018.

\bibitem{r13}
J. Dos Reis et al.,
``YOLO-Based Object Detection for Autonomous Navigation,''
\textit{IEEE Access},
vol.~7, pp.~124908--124920, 2019.

\bibitem{r14}
J. Jean and K. Lian,
``Robust Visual Servo Control Using Shape Parameters,''
\textit{IEEE/ASME Transactions on Mechatronics},
vol.~17, no.~1, pp.~30--40, 2012.

\bibitem{r15}
A. Zde\v{s}ar, D. Matko, and M. Bla\v{z}i\v{c},
``Visual Trajectory Tracking Control of Mobile Robots,''
\textit{IEEE Transactions on Control Systems Technology},
vol.~21, no.~6, pp.~2210--2217, 2013.

\bibitem{r16}
S. Mir and F. Filali,
``LTE and IEEE 802.11p for Vehicular Networking: A Performance Evaluation,''
\textit{EURASIP Journal on Wireless Communications and Networking},
2014.

\bibitem{r17}
K. Yousif, A. Bab-Hadiashar, and R. Hoseinnezhad,
``An Overview to Visual Odometry and Visual SLAM,''
\textit{Applied Bionics and Biomechanics},
2015.

\bibitem{r18}
C. Pati and R. Kala,
``Vision-Based Robot Following Using RGB-D Sensors,''
\textit{Procedia Computer Science},
vol.~115, pp.~586--593, 2017.

\bibitem{r19}
Z. Xu et al.,
``DSRC versus 4G-LTE for Connected Vehicle Applications: A Study on Field Experiments,''
\textit{Journal of Advanced Transportation},
2017.

\bibitem{r20}
Y. Zhang, Z. Zhu, and H. Wang,
``Global Homography Calibration for Vision-Based Robot Navigation,''
\textit{IEEE Transactions on Instrumentation and Measurement},
vol.~66, no.~10, pp.~2754--2763, 2017.

\bibitem{r21}
M. Fauadi et al.,
``A Review on Vision-Based Navigation Systems for Mobile Robots,''
\textit{International Journal of Advanced Robotic Systems},
2018.

\bibitem{r22}
J. M. Lozano Dom\'{i}nguez and T. J. Mateo Sanguino,
``Review on V2X, I2X, and P2X Communications and Their Applications,''
\textit{Sensors},
vol.~19, no.~12, p.~2756, 2019.

\bibitem{r23}
P. S. Santana et al.,
``Vehicle-Mounted Sensing for Mobile Air Quality Monitoring,''
\textit{IEEE Conference on Intelligent Transportation Systems},
2019.

\bibitem{r24}
P. Papadeas, G. Athanasopoulos, and K. Karantzalos,
``A Survey on Real-Time Semantic Segmentation for Autonomous Systems,''
\textit{IEEE Access},
vol.~9, pp.~54215--54240, 2021.

\bibitem{r25}
T. Petrov et al.,
``Performance Benchmark for DSRC and LTE-V2X Communication,''
\textit{Sensors},
vol.~21, no.~6, 2021.

\bibitem{r26}
A. Singh, R. Sharma, and P. Kumar,
``Machine Learning-Based Obstacle Avoidance and Path Tracking for Mobile Robots,''
\textit{IEEE Access},
vol.~10, pp.~112345--112356, 2022.

\bibitem{r27}
A. Saidi, H. Mekki, and M. Abid,
``Fuzzy-PID Control for Vision-Based Mobile Robot Navigation,''
\textit{Robotics and Autonomous Systems},
vol.~150, 2022.

\bibitem{r28}
C. Shin, E. Farag, H. Ryu, M. Zhou, and Y. Kim,
``Vehicle-to-Everything (V2X) Evolution from 4G to 5G in 3GPP,''
\textit{IEEE Access},
vol.~11, pp.~18689--18701, 2023.

\bibitem{r29}
J. Marroqu\'{i}n, R. Monroy, and L. Castillo,
``Embedded Computer Vision Navigation System for Mobile Robots,''
\textit{IEEE Access},
vol.~11, pp.~56789--56805, 2023.

\bibitem{r30}
M. M\i s\i r,
``DeepLab-Based Drivable Path Detection for Mobile Robot Navigation,''
\textit{IEEE Access},
vol.~12, pp.~12345--12360, 2024.

\bibitem{r31}
CVAT.ai Corporation,
\textit{CVAT: Computer Vision Annotation Tool},
2024.
[Online]. Available: \url{https://github.com/cvat-ai/cvat}

\bibitem{r32}
Ultralytics,
\textit{YOLOv8: Ultralytics YOLO Documentation},
Version 8.4.7, 2024.
[Online]. Available: \url{https://docs.ultralytics.com}

\bibitem{r33}
ONNX Runtime Development Team,
\textit{ONNX Runtime: Cross-Platform Inference Accelerator},
Version 1.23.2, 2024.
[Online]. Available: \url{https://onnxruntime.ai/docs/}

\bibitem{r34}
Raspberry Pi Foundation,
\textit{Raspberry Pi 5 Official Documentation},
2024.
[Online]. Available: \url{https://www.raspberrypi.com/documentation/}

\bibitem{r35}
Google Firebase,
\textit{Firebase Realtime Database Documentation},
Google LLC, 2024.
[Online]. Available: \url{https://firebase.google.com/docs/database}

\bibitem{r36}
MathWorks ThingSpeak,
\textit{ThingSpeak IoT Analytics Platform Documentation},
MathWorks, 2024.
[Online]. Available: \url{https://uk.mathworks.com/help/thingspeak/}

\bibitem{r37}
Pallets Projects,
\textit{Flask: A Lightweight WSGI Web Application Framework},
Version 3.x, 2024.
[Online]. Available: \url{https://flask.palletsprojects.com/}

\bibitem{r38}
Adafruit Industries,
\textit{Adafruit CircuitPython PCA9685 Library},
2024.
[Online]. Available: \url{https://github.com/adafruit/Adafruit_CircuitPython_PCA9685}

\bibitem{r39}
F. Pedregosa et al.,
``Scikit-learn: Machine Learning in Python,''
\textit{Journal of Machine Learning Research},
vol.~12, pp.~2825--2830, 2011.
[Online]. Available: \url{https://scikit-learn.org}

\bibitem{r40}
C. R. Harris et al.,
``Array Programming with NumPy,''
\textit{Nature},
vol.~585, pp.~357--362, 2020.
[Online]. Available: \url{https://numpy.org}

\bibitem{r41}
P. Virtanen et al.,
``SciPy 1.0: Fundamental Algorithms for Scientific Computing in Python,''
\textit{Nature Methods},
vol.~17, pp.~261--272, 2020.
[Online]. Available: \url{https://scipy.org}

\bibitem{r42}
J. D. Hunter,
``Matplotlib: A 2D Graphics Environment,''
\textit{Computing in Science \& Engineering},
vol.~9, no.~3, pp.~90--95, 2007.
[Online]. Available: \url{https://matplotlib.org}

\bibitem{r43}
OpenCV Development Team,
\textit{OpenCV: Open Source Computer Vision Library},
Version 4.x, 2024.
[Online]. Available: \url{https://opencv.org}

\bibitem{r44}
PySerial Development Team,
\textit{PySerial: Python Serial Port Extension},
Version 3.x, 2024.
[Online]. Available: \url{https://github.com/pyserial/pyserial}

\bibitem{r45}
Google Cloud,
\textit{Firebase Admin Python SDK},
Version 6.x, 2024.
[Online]. Available: \url{https://github.com/firebase/firebase-admin-python}

\bibitem{r46}
American Society of Heating, Refrigerating and Air-Conditioning Engineers (ASHRAE),
\textit{ASHRAE Standard 62.1: Ventilation for Acceptable Indoor Air Quality},
ASHRAE, Atlanta, GA, 2022.

\bibitem{r47}
World Health Organization (WHO),
\textit{WHO Global Air Quality Guidelines},
WHO, Geneva, 2021.
[Online]. Available: \url{https://www.who.int/publications/i/item/9789240034228}

\bibitem{r48}
O.\ Sepp\"{a}nen, W.\,J.\ Fisk, and M.\,J.\ Mendell,
``Association of Ventilation Rates and CO$_2$ Concentrations with Health
and Other Responses in Commercial and Institutional Buildings,''
\textit{Indoor Air},
vol.~9, no.~4, pp.~226--252, 1999.

\bibitem{r49}
U.\ Satish et al.,
``Is CO$_2$ an Indoor Pollutant? Direct Effects of Low-to-Moderate
CO$_2$ Concentrations on Human Decision-Making Performance,''
\textit{Environmental Health Perspectives},
vol.~120, no.~12, pp.~1671--1677, 2012.

\bibitem{r50}
Hanwei Electronics Co.,
\textit{MQ-135 Gas Sensor for Air Quality Datasheet},
Hanwei Electronics, Zhengzhou, China, 2013.
[Online]. Available: \url{https://www.hwsensor.com}

\end{thebibliography}

\chapter{Conclusion and Future Work}

This dissertation presented the design, development, and experimental
evaluation of a communication-integrated mobile environmental monitoring
framework implemented on a Raspberry Pi~5 and Raspberry Pi~3 embedded
computing platform pair. The proposed framework addressed the fundamental
limitation of conventional fixed-station environmental monitoring by
integrating an autonomous mobile robotic platform with real-time
environmental sensing and multi-channel cloud-enabled data communication
within a modular and scalable architecture.

The framework was realised through four tightly integrated components.
The web-based manual control system provided a validated hardware
foundation through remote teleoperation, confirming the correct
operation of all mechanical, sensing, and communication subsystems.
The YOLOv8n-seg path segmentation module, trained on a custom
289-image annotated dataset using the CVAT annotation platform,
achieved a mask mAP@50 of 0.988 and a mask mAP@50--95 of 0.736 on
the 57-image validation split. Systematic optimisation through ONNX
format conversion and 320$\times$320 resolution tuning achieved
real-time inference at 15.4\,FPS with 64.8\,ms latency on the
Raspberry Pi~5 CPU, representing a 1.57$\times$ speedup over the
PyTorch baseline and exceeding the minimum frame rate required for
stable closed-loop control.

The autonomous path-following subsystem, driven by a PID controller
with gains $K_p = 13.2$, $K_i = 0.0$, and $K_d = 0.475$, operated
through a four-thread concurrent architecture at execution rates of
20\,Hz for vision inference, 10\,Hz for obstacle monitoring, and
10\,Hz for motor command execution. A systematic experimental
evaluation comprising 20 manual and 20 autonomous trials on a
rectangular 152.4\,cm $\times$ 91.44\,cm test path demonstrated that
autonomous navigation achieved a mean RMSE of 2.119\,cm, representing
an 87.4\% reduction from the manual mean RMSE of 16.768\,cm, and a
tracking accuracy of 95.37\% compared to 63.32\% under manual
teleoperation. All three statistical hypothesis tests (Mann--Whitney U,
permutation, and independent $t$-test) unanimously confirmed the
improvement to be highly significant, with an exceptionally large
Cohen's $d = 8.944$ and a probability of superiority of 100\%.

The environmental monitoring subsystem acquired temperature, humidity,
and air quality data from an Arduino-interfaced DHT11 and MQ135 sensor
pair, applied linear calibration to estimate CO$_2$ concentration in
ppm, and implemented a six-level rule-based alert hierarchy aligned
with ASHRAE Standard~62.1 and WHO Indoor Air Quality Guidelines. All
acquired readings were simultaneously persisted to a local SQLite
database, pushed to a Firebase Realtime Database hosted at the
asia-southeast1 region, transmitted to a ThingSpeak IoT channel with
six mapped sensor fields, and served through a Flask live web dashboard
with gauge displays, trend charts, and a reading history table. All
four parallel output channels operated concurrently without mutual
blocking, confirming the effectiveness of the threading architecture.

The integrated two-processor architecture, with the Raspberry Pi~5
managing autonomous navigation and the Raspberry Pi~3 managing
environmental data communication, demonstrated effective isolation
between the real-time navigation workload and the cloud communication
workload. Both subsystems operated concurrently on the shared mobile
platform without observable interference. The framework contributes a
practical and validated approach to communication-integrated mobile
sensing that combines deep learning perception, real-time embedded
control, and multi-platform IoT data communication on affordable
embedded hardware. The modular design supports future extensions,
establishing a replicable foundation for advanced mobile environmental
monitoring and smart infrastructure applications.

% ── Future Work ───────────────────────────────────────────────
\section{Future Work}

Several directions for future research and development are identified
based on the limitations encountered during implementation and the
results of experimental evaluation.

\subsection{Real Wheel Odometry Integration}

The current trajectory logging and position estimation system uses
simulated dead-reckoning from vision deviation and speed commands,
which accumulates drift over time. Integration of physical wheel
encoders or odometry sensors would enable accurate measurement of
distance travelled, angular velocity, and absolute position, replacing
the simulated estimate with ground-truth trajectory data and
significantly improving performance evaluation reliability.

\subsection{Model Quantisation and Hardware Acceleration}

The current ONNX model at 320$\times$320 resolution achieves 15.4\,FPS
on the Raspberry Pi~5 CPU. INT8 quantisation of the ONNX model is
expected to provide a further 2$\times$ speedup. Hardware acceleration
using the Coral USB Accelerator or Hailo AI HAT for the Raspberry Pi~5
could reach 25 to 30\,FPS, enabling smoother control at higher robot
speeds without requiring a full GPU platform.

\subsection{Extended Environmental Sensing}

The current subsystem measures temperature, humidity, and air quality
through the DHT11 and MQ135 sensors. Future work could extend the
sensing suite to include a dedicated CO$_2$ sensor such as the MH-Z19
for direct calibrated measurement, a particulate matter sensor such as
the PMS5003 for PM2.5 and PM10 monitoring, a volatile organic compound
sensor, and noise level monitoring, providing a more comprehensive
environmental data profile for smart building and occupational health
applications.

\subsection{Multi-Robot Distributed Framework}

The current framework operates as a single mobile unit. A natural
extension is the deployment of multiple mobile monitoring units that
collaborate to cover larger spatial areas. Each unit would transmit
environmental data to a shared cloud backend, enabling spatial mapping
of environmental parameters. Distributed task allocation algorithms
could coordinate path assignments between units to maximise coverage
efficiency.

\subsection{Advanced Communication Integration}

Future work could replace the current Wi-Fi-based HTTP communication
with cellular LTE or 5G modules to enable monitoring in areas without
fixed network infrastructure. Integration of LTE-V2X communication
frameworks could further extend the system toward vehicular
environmental monitoring applications in outdoor and urban environments.

% ── References ───────────────────────────────────────────────
\newpage
\renewcommand{\bibname}{References}
\addcontentsline{toc}{chapter}{References}

\begin{thebibliography}{50}

\bibitem{r1}
R. Pushkaranand and M. Reddy,
``IoT-Based Vehicular Pollution Monitoring System,''
\textit{IEEE International Conference on Green Computing and Internet of Things (ICGCIoT)},
2015.

\bibitem{r2}
Y. Ben-Aboud, M. Ghogho, and A. Kobbane,
``A Research-Oriented Low-Cost Air Pollution Monitoring IoT Platform,''
\textit{IEEE Conference},
2020.

\bibitem{r3}
P. Papadimitratos et al.,
``Secure Vehicular Communication Systems: Design and Architecture,''
\textit{IEEE Communications Magazine},
vol.~46, no.~11, pp.~100--109, Nov. 2008.

\bibitem{r4}
G. Araniti, C. Campolo, M. Condoluci, A. Iera, and A. Molinaro,
``LTE for Vehicular Networking: A Survey,''
\textit{IEEE Communications Magazine},
vol.~51, no.~5, pp.~148--157, May 2013.

\bibitem{r5}
Y. Ma, J. Ko\v{s}eck\'{a}, and S. Sastry,
``Vision Guided Navigation for a Nonholonomic Mobile Robot,''
\textit{Proceedings of the IEEE International Conference on Robotics and Automation (ICRA)},
1999.

\bibitem{r6}
Y. Fang, W. Dixon, D. Dawson, and C. Abdallah,
``Homography-Based Visual Servoing of Wheeled Mobile Robots,''
\textit{IEEE Transactions on Robotics},
vol.~21, no.~3, pp.~406--414, 2005.

\bibitem{r7}
Y. Fang, W. Dixon, D. Dawson, and C. Abdallah,
``Homography-Based Visual Servo Regulation for Mobile Robots,''
\textit{IEEE Transactions on Robotics},
vol.~21, no.~3, pp.~406--414, 2005.

\bibitem{r8}
J. Chen, W. Dixon, D. Dawson, and M. McIntyre,
``Homography-Based Visual Servo Tracking Control of a Wheeled Mobile Robot,''
\textit{IEEE Transactions on Robotics},
2006.

\bibitem{r9}
E. Shaban and A. Nada,
``PID vs PIP Control for Mobile Robots,''
\textit{International Journal of Computer Applications},
vol.~117, no.~2, 2015.

\bibitem{r10}
L. Pacheco and N. Luo,
``PID vs MPC Control Strategies for Mobile Robots: A Comparative Study,''
\textit{IEEE Conference on Control Applications},
2015.

\bibitem{r11}
N. Noordin et al.,
``Comparison of PID and Fuzzy Logic Control for Mobile Robots,''
\textit{International Journal of Robotics Research},
2021.

\bibitem{r12}
M. Siam, S. Elkerdawy, M. Jagersand, and S. Yogamani,
``Deep Semantic Segmentation for Automated Driving: A Survey,''
\textit{IEEE Transactions on Intelligent Transportation Systems},
vol.~20, no.~2, pp.~641--659, 2018.

\bibitem{r13}
J. Dos Reis et al.,
``YOLO-Based Object Detection for Autonomous Navigation,''
\textit{IEEE Access},
vol.~7, pp.~124908--124920, 2019.

\bibitem{r14}
J. Jean and K. Lian,
``Robust Visual Servo Control Using Shape Parameters,''
\textit{IEEE/ASME Transactions on Mechatronics},
vol.~17, no.~1, pp.~30--40, 2012.

\bibitem{r15}
A. Zde\v{s}ar, D. Matko, and M. Bla\v{z}i\v{c},
``Visual Trajectory Tracking Control of Mobile Robots,''
\textit{IEEE Transactions on Control Systems Technology},
vol.~21, no.~6, pp.~2210--2217, 2013.

\bibitem{r16}
S. Mir and F. Filali,
``LTE and IEEE 802.11p for Vehicular Networking: A Performance Evaluation,''
\textit{EURASIP Journal on Wireless Communications and Networking},
2014.

\bibitem{r17}
K. Yousif, A. Bab-Hadiashar, and R. Hoseinnezhad,
``An Overview to Visual Odometry and Visual SLAM,''
\textit{Applied Bionics and Biomechanics},
2015.

\bibitem{r18}
C. Pati and R. Kala,
``Vision-Based Robot Following Using RGB-D Sensors,''
\textit{Procedia Computer Science},
vol.~115, pp.~586--593, 2017.

\bibitem{r19}
Z. Xu et al.,
``DSRC versus 4G-LTE for Connected Vehicle Applications: A Study on Field Experiments,''
\textit{Journal of Advanced Transportation},
2017.

\bibitem{r20}
Y. Zhang, Z. Zhu, and H. Wang,
``Global Homography Calibration for Vision-Based Robot Navigation,''
\textit{IEEE Transactions on Instrumentation and Measurement},
vol.~66, no.~10, pp.~2754--2763, 2017.

\bibitem{r21}
M. Fauadi et al.,
``A Review on Vision-Based Navigation Systems for Mobile Robots,''
\textit{International Journal of Advanced Robotic Systems},
2018.

\bibitem{r22}
J. M. Lozano Dom\'{i}nguez and T. J. Mateo Sanguino,
``Review on V2X, I2X, and P2X Communications and Their Applications,''
\textit{Sensors},
vol.~19, no.~12, p.~2756, 2019.

\bibitem{r23}
P. S. Santana et al.,
``Vehicle-Mounted Sensing for Mobile Air Quality Monitoring,''
\textit{IEEE Conference on Intelligent Transportation Systems},
2019.

\bibitem{r24}
P. Papadeas, G. Athanasopoulos, and K. Karantzalos,
``A Survey on Real-Time Semantic Segmentation for Autonomous Systems,''
\textit{IEEE Access},
vol.~9, pp.~54215--54240, 2021.

\bibitem{r25}
T. Petrov et al.,
``Performance Benchmark for DSRC and LTE-V2X Communication,''
\textit{Sensors},
vol.~21, no.~6, 2021.

\bibitem{r26}
A. Singh, R. Sharma, and P. Kumar,
``Machine Learning-Based Obstacle Avoidance and Path Tracking for Mobile Robots,''
\textit{IEEE Access},
vol.~10, pp.~112345--112356, 2022.

\bibitem{r27}
A. Saidi, H. Mekki, and M. Abid,
``Fuzzy-PID Control for Vision-Based Mobile Robot Navigation,''
\textit{Robotics and Autonomous Systems},
vol.~150, 2022.

\bibitem{r28}
C. Shin, E. Farag, H. Ryu, M. Zhou, and Y. Kim,
``Vehicle-to-Everything (V2X) Evolution from 4G to 5G in 3GPP,''
\textit{IEEE Access},
vol.~11, pp.~18689--18701, 2023.

\bibitem{r29}
J. Marroqu\'{i}n, R. Monroy, and L. Castillo,
``Embedded Computer Vision Navigation System for Mobile Robots,''
\textit{IEEE Access},
vol.~11, pp.~56789--56805, 2023.

\bibitem{r30}
M. M\i s\i r,
``DeepLab-Based Drivable Path Detection for Mobile Robot Navigation,''
\textit{IEEE Access},
vol.~12, pp.~12345--12360, 2024.

\bibitem{r31}
CVAT.ai Corporation,
\textit{CVAT: Computer Vision Annotation Tool},
2024.
[Online]. Available: \url{https://github.com/cvat-ai/cvat}

\bibitem{r32}
Ultralytics,
\textit{YOLOv8: Ultralytics YOLO Documentation},
Version 8.4.7, 2024.
[Online]. Available: \url{https://docs.ultralytics.com}

\bibitem{r33}
ONNX Runtime Development Team,
\textit{ONNX Runtime: Cross-Platform Inference Accelerator},
Version 1.23.2, 2024.
[Online]. Available: \url{https://onnxruntime.ai/docs/}

\bibitem{r34}
Raspberry Pi Foundation,
\textit{Raspberry Pi 5 Official Documentation},
2024.
[Online]. Available: \url{https://www.raspberrypi.com/documentation/}

\bibitem{r35}
Google Firebase,
\textit{Firebase Realtime Database Documentation},
Google LLC, 2024.
[Online]. Available: \url{https://firebase.google.com/docs/database}

\bibitem{r36}
MathWorks ThingSpeak,
\textit{ThingSpeak IoT Analytics Platform Documentation},
MathWorks, 2024.
[Online]. Available: \url{https://uk.mathworks.com/help/thingspeak/}

\bibitem{r37}
Pallets Projects,
\textit{Flask: A Lightweight WSGI Web Application Framework},
Version 3.x, 2024.
[Online]. Available: \url{https://flask.palletsprojects.com/}

\bibitem{r38}
Adafruit Industries,
\textit{Adafruit CircuitPython PCA9685 Library},
2024.
[Online]. Available: \url{https://github.com/adafruit/Adafruit_CircuitPython_PCA9685}

\bibitem{r39}
F. Pedregosa et al.,
``Scikit-learn: Machine Learning in Python,''
\textit{Journal of Machine Learning Research},
vol.~12, pp.~2825--2830, 2011.
[Online]. Available: \url{https://scikit-learn.org}

\bibitem{r40}
C. R. Harris et al.,
``Array Programming with NumPy,''
\textit{Nature},
vol.~585, pp.~357--362, 2020.
[Online]. Available: \url{https://numpy.org}

\bibitem{r41}
P. Virtanen et al.,
``SciPy 1.0: Fundamental Algorithms for Scientific Computing in Python,''
\textit{Nature Methods},
vol.~17, pp.~261--272, 2020.
[Online]. Available: \url{https://scipy.org}

\bibitem{r42}
J. D. Hunter,
``Matplotlib: A 2D Graphics Environment,''
\textit{Computing in Science \& Engineering},
vol.~9, no.~3, pp.~90--95, 2007.
[Online]. Available: \url{https://matplotlib.org}

\bibitem{r43}
OpenCV Development Team,
\textit{OpenCV: Open Source Computer Vision Library},
Version 4.x, 2024.
[Online]. Available: \url{https://opencv.org}

\bibitem{r44}
PySerial Development Team,
\textit{PySerial: Python Serial Port Extension},
Version 3.x, 2024.
[Online]. Available: \url{https://github.com/pyserial/pyserial}

\bibitem{r45}
Google Cloud,
\textit{Firebase Admin Python SDK},
Version 6.x, 2024.
[Online]. Available: \url{https://github.com/firebase/firebase-admin-python}

\bibitem{r46}
American Society of Heating, Refrigerating and Air-Conditioning Engineers (ASHRAE),
\textit{ASHRAE Standard 62.1: Ventilation for Acceptable Indoor Air Quality},
ASHRAE, Atlanta, GA, 2022.

\bibitem{r47}
World Health Organization (WHO),
\textit{WHO Global Air Quality Guidelines},
WHO, Geneva, 2021.
[Online]. Available: \url{https://www.who.int/publications/i/item/9789240034228}

\bibitem{r48}
O.\ Sepp\"{a}nen, W.\,J.\ Fisk, and M.\,J.\ Mendell,
``Association of Ventilation Rates and CO$_2$ Concentrations with Health
and Other Responses in Commercial and Institutional Buildings,''
\textit{Indoor Air},
vol.~9, no.~4, pp.~226--252, 1999.

\bibitem{r49}
U.\ Satish et al.,
``Is CO$_2$ an Indoor Pollutant? Direct Effects of Low-to-Moderate
CO$_2$ Concentrations on Human Decision-Making Performance,''
\textit{Environmental Health Perspectives},
vol.~120, no.~12, pp.~1671--1677, 2012.

\bibitem{r50}
Hanwei Electronics Co.,
\textit{MQ-135 Gas Sensor for Air Quality Datasheet},
Hanwei Electronics, Zhengzhou, China, 2013.
[Online]. Available: \url{https://www.hwsensor.com}

\end{thebibliography}